\documentclass[10pt,a4paper]{article}
\usepackage[margin=2.0cm]{geometry}
\usepackage{fontspec}
\usepackage{amsmath,amssymb,amsfonts,amsthm,mathtools}
\usepackage{booktabs,tabularx,colortbl,array,multirow}
\usepackage[table]{xcolor}
\usepackage{fancyhdr}
\usepackage{titlesec}
\usepackage{microtype}
\usepackage{hyperref}

\newfontfamily\thaifont{Thonburi}

\hypersetup{
    colorlinks=true,
    linkcolor=blue!70!black,
    citecolor=blue!70!black,
    urlcolor=blue!70!black
}

\pagestyle{fancy}
\fancyhf{}
\rhead{\footnotesize \textbf{NEXVIA AEGIS V9-M4 H3 CI \& CRN Claim Audit}}
\lhead{\footnotesize Authoritative Protocol Verification, Method Reconciliation \& Evidentiary Audit}
\rfoot{\footnotesize Page \thepage}
\lfoot{\footnotesize Document ID: AEGIS-V9-M4-CRN-18}

\titleformat{\section}{\large\bfseries}{\thesection}{1em}{}
\titleformat{\subsection}{\normalsize\bfseries}{\thesubsection}{1em}{}

\begin{document}

\begin{center}
    {\LARGE \textbf{NEXVIA AEGIS V9-M4 Methodological Audit}}\\[6pt]
    {\Large \textbf{Final Hypothesis H3 Confidence-Bound Implementation Verification}}\\[3pt]
    {\Large \textbf{and Common-Random-Numbers Variance-Reduction Claim Audit}}\\[10pt]
    \textbf{Autonomous Verification Subsystem (Google DeepMind Antigravity)}\\
    \textit{Git Branch: v9.0-runtime-assurance \quad $\cdot$ \quad Reported Prior Commit: 0d8350c \quad $\cdot$ \quad Date: October 25, 2026}\\[12pt]
\end{center}

\noindent\fbox{
\parbox{0.97\textwidth}{
\textbf{CRITICAL NON-AUTHORIZATION DIRECTIVE:}\\
This audit document resolves two specific pre-execution questions: (1) Does the implemented Hypothesis $H_3$ noninferiority confidence bound match the authoritative registered experiment protocol? (2) Are claims regarding variance reduction from the Common Random Numbers (CRN) design supported by empirical and theoretical evidence? \textbf{This document does NOT constitute authorization to execute Milestone V9-M4.} All evaluations in this report were conducted strictly using synthetic test fixtures, in-memory models, and mock environments. Production holdout manifests (\texttt{test\_mix\_manifest\_r1.json} and \texttt{test\_scen\_manifest\_r1.json}) and output directories (\texttt{results/v9/test\_mix/} and \texttt{results/v9/test\_scen/}) do not exist. Reserved holdout seed blocks ($73000..74439$ for mixture holdout and $75000..75119$ for scenario holdout) remain completely uninstantiated and untouched. Live execution requires explicit authorization from the project owner.
}
}

\vspace{10pt}

\begin{abstract}
\noindent This audit resolves two foundational methodological questions prior to production holdout authorization: (1) Does the implemented Hypothesis $H_3$ True-Success Rate (TSR) noninferiority confidence bound conform to the authoritative preregistered protocol? (2) Are claims regarding variance reduction attributable to the Common Random Numbers (CRN) design supported by available evidence?

First, comparing governing protocol texts (\texttt{docs/v9\_experiment\_protocol.md}) with prior reports, we reconcile the protocol's dual parenthetical reference to \textit{``Newcombe/Wald paired difference CI''} with the pre-execution report's specification of a paired-difference Student's $t$ lower confidence bound ($\mathrm{LCB}_{95}$). Mathematical analysis and synthetic boundary testing on the registered sample size $N = 1,200$ establish that the sample standard error $\frac{s_d}{\sqrt{N}}$ is the finite-sample Student's $t$ formulation of the paired Wald interval. At $N = 1,200$, the Student's $t$ LCB and Newcombe's profile-score LCB (Method 10, two-sided 90\% confidence / one-sided 95\% level) agree to within $\pm 0.00002$ ($0.002$ percentage points) across realistic discordant rates, yielding 100\% decision concordance across all boundary cases. To ensure absolute protocol fidelity and eliminate ambiguity, we updated \texttt{experiments/v9\_evaluate\_holdout.py} to compute and report both bounds (\texttt{lcb95\_delta\_tsr} and \texttt{newcombe\_lcb95\_delta\_tsr}) alongside the complete $2 \times 2$ contingency table.

Second, auditing CRN variance-reduction claims, we distinguish structural pairing mechanisms from empirical outcome distributions. While the CRN architecture enforces matched initial conditions ($t=0$) and shared scenario seeds, variance reduction $\operatorname{Var}(\overline{\Delta\mathrm{TSR}}) < \operatorname{Var}_{\text{unpaired}}$ mathematically requires positive covariance $\operatorname{Cov}(Y_{\pi}, Y_{\text{base}}) > 0$. On historical benchmark data ($N = 500$), empirical correlation was strongly positive ($\rho = 0.5931$), producing an empirical standard error reduction of $36.21\%$. However, on the unexecuted holdout partition ($\mathcal{D}_{\text{test}}^{\text{scen}}$), holdout data has not been generated or simulated. Consequently, variance reduction cannot be claimed as an established holdout fact or an unconditional architectural guarantee; rather, the CRN design \textit{enables} paired evaluation and is \textit{intended to exploit positive correlation}, but the actual reduction in uncertainty is an empirical property to be measured post-authorization.

Four dedicated synthetic regression tests (\texttt{test\_56}--\texttt{test\_59}) verify numerical convergence, boundary behavior, dual-bound reporting, and the empirical contingency of variance reduction. All 59 preflight tests and all 170 full repository tests pass with zero failures.
\end{abstract}

\vspace{6pt}
\noindent \textbf{Keywords:} Hypothesis H3, Noninferiority Confidence Bound, Student's $t$ Wald LCB, Newcombe Paired Score Interval, Common Random Numbers (CRN), Covariance Structure, Holdout Preservation.

\section{Executive Summary}

\subsection{Executive Summary (English)}
\begin{enumerate}
    \item \textbf{Confidence Bound Reconciliation:} Protocol Section 3.1.1 refers parenthetically to \textit{``Newcombe/Wald paired difference CI''}, while Section 3.2.4 derives power from the Wald standard error $\mathrm{SE}_0 = \sqrt{q/N}$, and pre-execution reports specify a paired Student's $t$ bound. We prove that at $N = 1,200$, the finite-sample Student's $t$ Wald LCB and Newcombe's paired score LCB (Method 10) are virtually indistinguishable ($|\Delta| \le 0.00002$) and yield identical noninferiority decisions.
    \item \textbf{Dual-Bound Implementation:} Updated \texttt{evaluate\_scenario\_holdout} in \texttt{experiments/v9\_evaluate\_holdout.py} to calculate and report both \texttt{lcb95\_delta\_tsr} (Student's $t$ Wald) and \texttt{newcombe\_lcb95\_delta\_tsr} (Newcombe 10) for both V8.1 Martingale and V7 MPC, guaranteeing compliance under either interpretation.
    \item \textbf{CRN Variance-Reduction Evidentiary Audit:} We formally qualify prior claims that CRN ``guarantees'' variance reduction. Mathematically, pairing reduces standard error if and only if $\operatorname{Cov}(Y_{\pi}, Y_{\text{base}}) > 0$. While historical data confirms $\rho \approx 0.593$ (yielding 36.2\% SE reduction between V8 and V7), the correlation between V9 and baselines on holdout data remains an unmeasured empirical property until execution.
    \item \textbf{Independent Synthetic Verification:} Verified all boundary cases on $N = 1,200$ matched pairs: all agreeing (100\% and 50\%), balanced discordant ($q = 0.204$), slight deficit ($\Delta = -0.015$), near boundary ($\Delta = -0.028333$ passes, $\Delta = -0.030$ fails), and nominal TSR failure ($112/120 = 0.9333 < 0.94$ fails).
    \item \textbf{Synthetic Test Suite:} Added unit tests \texttt{test\_56} through \texttt{test\_59} in \texttt{tests/test\_v9\_holdout\_preflight.py}. Preflight tests achieve 59/59 passing (1.90s); full regression suite achieves 170/170 passing (28.16s).
    \item \textbf{Holdout Independence:} Zero production holdout manifests, datasets, or RNG streams were created or accessed. The frozen policy hash (\texttt{5a89d606...}) and calibration artifacts remain 100\% identical.
\end{enumerate}

\subsection{\thaifont บทสรุปสำหรับผู้บริหาร (Executive Summary in Thai)}
\begin{enumerate}
    \item \textbf{\thaifont การพิสูจน์ความสอดคล้องของช่วงความเชื่อมั่น H3:} {\thaifont เอกสารพิธีสารหมวด 3.1.1 ระบุในวงเล็บถึง ``Newcombe/Wald paired difference CI'' ในขณะที่หมวด 3.2.4 คำนวณกำลังการทดสอบจาก Wald Standard Error และรายงานก่อนหน้าระบุเป็น Student's $t$ Lower Confidence Bound ($\mathrm{LCB}_{95}$) การพิสูจน์ทางคณิตศาสตร์และชุดข้อมูลสังเคราะห์ขนาด $N = 1,200$ ยืนยันว่าวิธี Student's $t$ Wald และวิธี Newcombe Method 10 ให้ผลลัพธ์ใกล้เคียงกันอย่างยิ่ง (ผลต่างไม่เกิน $0.00002$ หรือ $0.002\%$) และให้ผลการตัดสินความไม่ด้อยกว่าเหมือนกันทุกกรณี}
    \item \textbf{\thaifont การรายงานช่วงความเชื่อมั่นคู่ขนาน:} {\thaifont ปรับปรุงฟังก์ชัน \texttt{evaluate\_scenario\_holdout} ใน \texttt{experiments/v9\_evaluate\_holdout.py} ให้คำนวณและรายงานทั้ง \texttt{lcb95\_delta\_tsr} (Student's $t$) และ \texttt{newcombe\_lcb95\_delta\_tsr} (Newcombe Score) พร้อมตารางนับ $2 \times 2$ สำหรับทั้ง V8.1 และ V7 เพื่อขจัดความคลุมเครือและคงความสอดคล้องกับข้อกำหนดทุกประการ}
    \item \textbf{\thaifont การตรวจสอบข้ออ้างการลดความแปรปรวนของ CRN:} {\thaifont ทำการแก้ไขข้อความในรายงานก่อนหน้าที่ระบุว่า CRN ``รับประกัน'' การลดความแปรปรวน โดยในทางสถิติ การจับคู่จะลด Standard Error ได้ก็ต่อเมื่อผลลัพธ์มีความสัมพันธ์เชิงบวก ($\operatorname{Cov} > 0$) เท่านั้น แม้ข้อมูลประวัติศาสตร์จะยืนยันค่า $\rho \approx 0.593$ (ลด SE ได้ 36.2\%) แต่บนชุดข้อมูล Holdout จริงที่ยังไม่เคยรัน ค่าความสัมพันธ์นี้ยังถือเป็นตัวแปรเชิงประจักษ์ที่ต้องรอการวัดผลหลังได้รับอนุมัติ}
    \item \textbf{\thaifont การคำนวณและทดสอบกรณีขอบเขต:} {\thaifont ทดสอบกรณีสังเคราะห์อย่างอิสระบนกลุ่มตัวอย่าง $N = 1,200$ ครบถ้วนทุกมิติ: ผลลัพธ์ตรงกันทั้งหมด, ผลลัพธ์ขัดแย้งแบบสมดุล, อัตราสำเร็จด้อยกว่าเล็กน้อย, ขอบเขตวิกฤตความไม่ด้อยกว่า ($\Delta = -0.028333$ ผ่าน, $\Delta = -0.030$ ไม่ผ่าน) และกรณี TSR สภาวะปกติไม่ถึงเกณฑ์ ($112/120 < 0.94$)}
    \item \textbf{\thaifont การขยายชุดทดสอบเชิงสังเคราะห์:} {\thaifont เพิ่มการทดสอบ \texttt{test\_56} ถึง \texttt{test\_59} ส่งผลให้ชุดทดสอบ Preflight ผ่าน 59/59 รายการ (1.90 วินาที) และชุดทดสอบระบบรวมผ่าน 170/170 รายการ (28.16 วินาที)}
    \item \textbf{\thaifont การคงความบริสุทธิ์ของข้อมูล Holdout:} {\thaifont ยืนยันว่าไม่มีการสร้างไฟล์ Manifest หรือรันซิมูเลชันบนข้อมูล Holdout จริง และค่าแฮชของนโยบายควบคุมยังคงเดิม 100\%}
\end{enumerate}

\section{Phase 1 — Audit of Registered H3 Confidence-Bound Method}

\subsection{Cross-Document Reconciliation of Protocol Specifications}
A rigorous textual and mathematical comparison was conducted across all governing documents:

\begin{table}[h!]
\centering
\small
\begin{tabularx}{\textwidth}{l p{3.5cm} X p{3.2cm}}
\toprule
\textbf{Document Source} & \textbf{Section \& Line} & \textbf{Exact Phrasing} & \textbf{Method Terminology} \\
\midrule
\texttt{v9\_experiment\_protocol.md} & §3.1.1 (line 30) & \textit{``One-sided 95\% paired lower confidence bound (Newcombe/Wald paired difference CI on the 1,200 matched \texttt{(scenario, seed)} pairs) strictly exceeds -0.05...''} & Dual designation: ``Newcombe/Wald paired difference CI'' \\
\texttt{v9\_experiment\_protocol.md} & §3.2.4 (lines 112--118) & \textit{``paired standard error under null equality is $\mathrm{SE}_0 = \sqrt{q / N}$, two-sided 95\% CI half-width is $1.96\,\mathrm{SE}_0$, and MDE is $(z_{1-\alpha_{\mathrm{NI}}} + z_{1-\beta})\sqrt{q / N}$''} & Wald asymptotic paired difference formulation \\
\texttt{v9\_experiment\_protocol.md} & Table 3.5.1 (line 247) & \textit{``1-sided 95\% paired LCB on $\Delta\mathrm{TSR} = \mathrm{TSR}(\hat{\pi}) - \mathrm{TSR}(\pi_{\mathrm{ref}})$ across $N=1,200$ matched $\mathcal{S}_{10}$ pairs...''} & 1-sided 95\% paired LCB at margin $-0.05$ \\
Report 04 & §2.2 vs §3.1 & §2.2 applies Newcombe Method 10 to $H_2$ ($\Delta R_{\mathrm{inc}}$). §3.1 applies Student's $t$ paired LCB $\bar{d} - t_{0.95, N-1} \frac{s_d}{\sqrt{N}}$ to $H_3$ ($\Delta\mathrm{TSR}$). & Newcombe for $H_2$; Student's $t$ LCB for $H_3$ \\
Report 13 & §3.2 & Evaluated paired difference Student's $t$ lower bound for $H_3$. & Student's $t$ paired LCB \\
Report 16 \& 17 & §2.1 \& §4.1 & Confirmed V8.1 primary and V7 co-evaluation under paired LCB $> -0.05$ on $S_{10}$. & Paired LCB $> -0.05$ \\
\bottomrule
\end{tabularx}
\caption{Reconciliation of statistical confidence bound terminology across governing repository documents.}
\end{table}

\subsection{Mathematical Formulation of the Estimators}
The parenthetical phrase in Section 3.1.1 combines two standard biometric procedures for paired binary differences:

\paragraph{1. Finite-Sample Wald Student's $t$ Paired Lower Bound:}
Let $Y_{1, i}, Y_{2, i} \in \{0, 1\}$ represent binary outcomes on matched pair $i \in \{1, \dots, N\}$. The paired difference is $D_i = Y_{1, i} - Y_{2, i} \in \{-1, 0, 1\}$. The sample mean difference is $\overline{D} = \frac{n_{10} - n_{01}}{N} = \widehat{\mathrm{TSR}}_1 - \widehat{\mathrm{TSR}}_2$. The sample variance of differences is:
\begin{equation}
s_D^2 = \frac{1}{N-1}\sum_{i=1}^N (D_i - \overline{D})^2 = \frac{N}{N-1}\left( \frac{n_{10} + n_{01}}{N} - \overline{D}^2 \right)
\end{equation}
The standard error of the mean difference is $\mathrm{SE}(\overline{D}) = \frac{s_D}{\sqrt{N}} = \sqrt{\frac{p_{10} + p_{01} - (p_{10} - p_{01})^2}{N-1}}$. Substituting the Student's $t$ critical value $t_{1-\alpha, N-1}$ yields:
\begin{equation}
\mathrm{LCB}_{95}^{\text{Wald-}t} = \overline{D} - t_{0.95, N-1} \frac{s_D}{\sqrt{N}}
\end{equation}
This is the formulation implemented in \texttt{paired\_difference\_lcb}.

\paragraph{2. Newcombe Paired Difference Interval (Method 10):}
Newcombe (1998, Method 10) forms the paired score interval from the individual Wilson score confidence limits $[l_1, u_1]$ and $[l_2, u_2]$ for marginal proportions $p_1 = (n_{11} + n_{10})/N$ and $p_2 = (n_{11} + n_{01})/N$, coupled by the sample phi-coefficient correlation $\hat{\phi} = \frac{n_{11} n_{00} - n_{10} n_{01}}{\sqrt{(n_{11} + n_{10})(n_{01} + n_{00})(n_{11} + n_{01})(n_{10} + n_{00})}}$:
\begin{equation}
\mathrm{LCB}_{1-\alpha}^{\text{Newcombe}} = (p_1 - p_2) - \sqrt{(p_1 - l_1)^2 - 2\hat{\phi}(p_1 - l_1)(u_2 - p_2) + (u_2 - p_2)^2}
\end{equation}
where $l_1$ and $u_2$ are the Wilson limits at two-sided confidence $1 - 2\alpha = 0.90$.

\section{Phase 2 — Mathematical Comparison and Synthetic Boundary Verification}

\subsection{Numerical Comparison on Registered Sample Size ($N = 1,200$)}
We independently computed both estimators across key synthetic boundary fixtures on $N = 1,200$:

\begin{table}[h!]
\centering
\small
\begin{tabularx}{\textwidth}{c l c c c c c c c c}
\toprule
\textbf{\#} & \textbf{Case Description} & $n_{11}$ & $n_{10}$ & $n_{01}$ & $n_{00}$ & $\Delta\mathrm{TSR}$ & \textbf{Student's $t$ LCB} & \textbf{Newcombe LCB} & \textbf{Decision} \\
\midrule
1 & All pairs agreeing (100\% TS) & 1,200 & 0 & 0 & 0 & $+0.0000$ & $0.000000$ & $-0.002250$ & \textbf{PASS} (both) \\
2 & All pairs agreeing (50\% TS) & 600 & 0 & 0 & 600 & $+0.0000$ & $0.000000$ & $0.000000$ & \textbf{PASS} (both) \\
3 & Balanced discordant ($q = 0.204$) & 800 & 122 & 122 & 156 & $+0.0000$ & $-0.021437$ & $-0.021421$ & \textbf{PASS} (both) \\
4 & Slight deficit ($\Delta = -0.015$) & 800 & 113 & 131 & 156 & $-0.0150$ & $-0.036425$ & $-0.036410$ & \textbf{PASS} (both) \\
5 & Near boundary (Pass: $\Delta = -0.0283$) & 800 & 105 & 139 & 156 & $-0.0283$ & $-0.049728$ & $-0.049713$ & \textbf{PASS} (both) \\
6 & Below boundary (Fail: $\Delta = -0.030$) & 800 & 104 & 140 & 156 & $-0.0300$ & $-0.051389$ & $-0.051374$ & \textbf{FAIL} (both) \\
7 & Nominal TSR failure ($\widehat{\mathrm{TSR}}_{\mathrm{nom}} < 0.94$) & 800 & 122 & 122 & 156 & $+0.0000$ & $-0.021437$ & $-0.021421$ & \textbf{FAIL} (nom) \\
\bottomrule
\end{tabularx}
\caption{Numerical comparison of Student's $t$ Wald LCB vs. Newcombe Method 10 LCB on $N = 1,200$ matched pairs.}
\end{table}

Across all non-degenerate cases ($q > 0$), the difference between the Student's $t$ Wald LCB and Newcombe LCB is less than $0.00002$ ($0.002\%$). Both methods achieve 100\% decision concordance. To eliminate ambiguity, \texttt{evaluate\_scenario\_holdout} in \texttt{experiments/v9\_evaluate\_holdout.py} now outputs both bounds: \texttt{lcb95\_delta\_tsr} and \texttt{newcombe\_lcb95\_delta\_tsr}.

\section{Phase 3 — Audit of Common Random Numbers (CRN) Variance-Reduction Claims}

\subsection{Decomposition of Architectural Guarantees vs. Empirical Outcomes}
Prior Report 17 marked Guarantee G5 (Paired Variance Reduction) as ``YES: Guaranteed''. We formally decompose the CRN properties into distinct technical levels:

\begin{table}[h!]
\centering
\small
\begin{tabularx}{\textwidth}{l X c p{3.2cm}}
\toprule
\textbf{Property} & \textbf{Technical Description} & \textbf{Status} & \textbf{Evidentiary Basis} \\
\midrule
\textbf{1. Paired Identity} & Episodes across policies matched by \texttt{(scenario, seed)}. & Enforced & Key equality assertion. \\
\textbf{2. Matched Initial Conditions} & Target position $\mathbf{p}_{\mathrm{target}}$ and sensor noise at $t=0$ bitwise identical. & Enforced & Proven in \texttt{test\_49}. \\
\textbf{3. Separate Mutable States} & Independent \texttt{Generator} instances spawned via \texttt{SeedSequence}. & Enforced & Proven in \texttt{test\_48}, \texttt{50}. \\
\textbf{4. Timestep Noise Alignment} & Noise draws match at $t > 0$ across controllers. & Qualified & Streams offset upon anchor queries (Model A). \\
\textbf{5. Outcome Covariance} & Positive correlation $\operatorname{Cov}(Y_1, Y_2) > 0$ across paired outcomes. & Empirical & $\rho = 0.593$ on benchmark; unmeasured on holdout. \\
\textbf{6. Variance Reduction} & $\operatorname{Var}(\overline{\Delta\mathrm{TSR}}) < \operatorname{Var}_{\text{unpaired}}$ achieved by pairing. & Empirical & Holds iff $\operatorname{Cov} > 0$. Not an architectural guarantee. \\
\bottomrule
\end{tabularx}
\caption{Decomposition of CRN random-number properties into architectural mechanisms and empirical outcomes.}
\end{table}

\subsection{Historical Evidence and Prospective Holdout Status}
For paired binary outcomes across $N$ matched episodes:
\begin{equation}
\operatorname{Var}(\overline{D}) = \frac{\operatorname{Var}(Y_1) + \operatorname{Var}(Y_2) - 2\operatorname{Cov}(Y_1, Y_2)}{N} < \frac{\operatorname{Var}(Y_1) + \operatorname{Var}(Y_2)}{N} \iff \operatorname{Cov}(Y_1, Y_2) > 0
\end{equation}
On historical benchmark data ($N = 500$, seeds $60000..60049$ on $\mathcal{S}_{10}$), comparing V8.1 and V7 yields:
\begin{equation}
n_{11} = 208, \quad n_{10} = 43, \quad n_{01} = 59, \quad n_{00} = 190, \quad q = 0.2040
\end{equation}
\begin{equation}
\operatorname{Cov}(Y_{\mathrm{V8}}, Y_{\mathrm{V7}}) = +0.14793, \quad \rho = +0.5931, \quad \mathrm{SE}_{\text{unpaired}} = 0.03159, \quad \mathrm{SE}_{\text{paired}} = 0.02015 \quad (\text{SE reduction} = 36.21\%)
\end{equation}
However, on the unexecuted holdout partition ($\mathcal{D}_{\text{test}}^{\text{scen}}$), holdout data has never been generated or simulated. Consequently, variance reduction cannot be claimed as an established holdout fact or an unconditional architectural guarantee. Prior wording stating G5 is ``Guaranteed: YES'' is formally qualified: the CRN design \textit{enables} paired comparisons and is \textit{structurally intended to exploit positive correlation}, but the actual variance reduction is an empirical property that will be measured post-authorization.

\section{Synthetic Regression Testing and Verification}
Four synthetic tests (\texttt{test\_56} through \texttt{test\_59}) were added to \texttt{tests/test\_v9\_holdout\_preflight.py}:
\begin{itemize}
    \item \texttt{test\_56}: Verifies \texttt{newcombe\_paired\_difference\_lcb} across all 7 boundary fixtures.
    \item \texttt{test\_57}: Proves 100\% decision concordance and $|\Delta| \le 0.00005$ across discordant grids.
    \item \texttt{test\_58}: Confirms that \texttt{evaluate\_scenario\_holdout} outputs both LCB metrics and contingency tables.
    \item \texttt{test\_59}: Validates the covariance requirement $\operatorname{Var}_{\text{paired}} < \operatorname{Var}_{\text{unpaired}} \iff \operatorname{Cov} > 0$ across correlated, independent, and anti-correlated synthetic data.
\end{itemize}

\paragraph{Recorded Execution Results:}
\begin{itemize}
    \item \texttt{PYTHONPATH=src uv run pytest -v tests/test\_v9\_holdout\_preflight.py}: \textbf{59 passed in 1.90s (Exit Code: 0)}.
    \item \texttt{PYTHONPATH=src uv run pytest -q}: \textbf{170 passed in 28.16s (Exit Code: 0)}.
\end{itemize}

\section{Artifact Provenance and Checksums}

\begin{table}[h!]
\centering
\small
\begin{tabularx}{\textwidth}{l X c}
\toprule
\textbf{Artifact Path} & \textbf{SHA-256 Checksum} & \textbf{Status} \\
\midrule
\texttt{results/v9/cal/selected\_policy.json} & \texttt{1c316d2d97ccf1a93f2dcd4c4196aefbb3817f7bc951dee5b0290468f7de7b87} & UNCHANGED \\
\texttt{results/v9/cal/calibration\_report.json} & \texttt{18ab85ab7d54625e4bb0d65de4786fe2896357768ed562cbfdee966a5d133337} & UNCHANGED \\
\texttt{experiments/v9\_evaluate\_holdout.py} & \textit{Working tree (dual LCB reporting implemented)} & UPDATED \\
\texttt{tests/test\_v9\_holdout\_preflight.py} & \textit{Working tree (tests 56--59 added)} & UPDATED \\
\bottomrule
\end{tabularx}
\caption{Integrity verification of frozen policy, calibration artifacts, and updated test suite.}
\end{table}

\section{Conclusion and Recommendation}
\begin{center}
    \fbox{\textbf{FINAL STATUS: H3\_CI\_CRN\_CLAIMS\_CONFIRMED\_WITH\_QUALIFICATIONS}}
\end{center}
The H3 noninferiority confidence bound conforms to registered protocol specifications under both Student's $t$ Wald and Newcombe profile-score formulations ($|\Delta| \le 0.00002$, 100\% decision concordance). Dual reporting is now implemented. CRN variance reduction is mathematically clarified as an enabled empirical property verified historically ($\rho = 0.593$, SE reduction 36.2\%) and properly qualified on prospective holdout data.

\begin{thebibliography}{9}
\bibitem{newcombe1998} R. G. Newcombe, ``Improved confidence intervals for the difference between binomial proportions based on paired data,'' \textit{Statistics in Medicine}, vol. 17, no. 22, pp. 2635--2650, 1998.
\bibitem{wald1943} A. Wald, ``Tests of statistical hypotheses concerning several parameters when the number of observations is large,'' \textit{Transactions of the American Mathematical Society}, vol. 54, no. 3, pp. 426--482, 1943.
\end{thebibliography}

\end{document}
